\section{Benchmark Details}
\label{sec:benchmark-details}
% 【中文】本文件已含 v1.7 时期中文对照注释（`% <chinese>` 行内格式，功能等价；2026-11-18 中文对照令核验保留，未改以守判据③纯新增）。

% 附录 A：benchmark 构建细节（参考 SWE-bench Appendix A "Benchmark Details"）
% 附录 A：基准构建细节（参考 SWE-bench Appendix A "Benchmark Details"）

\subsection{Task Example}
\label{sec:task-example}
% \subsection{任务示例}

[TODO: show one complete task: prompt + assets + check.py + test\_rationale. Suggested: a medium-difficulty schematic\_netlist case that demonstrates the full design\_flow contract.]
% [待办：展示一个完整任务：提示词 + 附件 + check.py + test\_rationale。建议选一个中等难度的 schematic\_netlist 案例，以演示完整的 design\_flow 契约。]

\subsection{Scan Pipeline Details}
\label{sec:scan-pipeline}
% \subsection{扫描流程细节}

We surveyed the engineering workstation fleet and scanned active production machines, yielding several thousand sessions with EDA keyword matches (mc $\geq 15$). Sessions are ranked by match-count (mc)---the number of lines hitting any EDA keyword. The mc $\geq 30$ threshold yields the high-quality candidate pool from which the \emph{Full} set (\FullScale{} tasks, continuously growing, count not reported in the main paper) is curated; a \emph{Lite} subset of \LiteSet{} tasks is filtered from Full by domain-coverage balance (all 13 domains preserved) and chain-length distribution (approximating Full's) for general evaluation under bounded token/time budgets. The exact workstation count, user count, and raw session inventory are company-internal and omitted here; the curation methodology is fully described so that another organization with comparable trace access can replicate the pipeline.
% 我们调研了工程工作站集群并扫描了生产环境中的活跃机器，产出数千个含 EDA 关键词命中的会话（mc $\geq 15$）。会话按匹配计数（mc）——命中任意 EDA 关键词的行数——排序。mc $\geq 30$ 的阈值产出高质量候选池，\emph{Full} 集（\FullScale{} 个任务、持续增长、正文不报告具体数字）即从中整理；\emph{Lite} 子集 \LiteSet{} 个任务则从 Full 按领域覆盖平衡（保留全部 13 域）与链长分布（近似 Full）筛选，用于在有限 token/时间预算下的一般评测。确切的工作站数、用户数与原始会话清单属公司内部信息，此处略去；整理方法论已完整描述，使另一家具有可比 trace 访问权限的组织可以复现该流程。

\textbf{Two-stage prefilter.} A rule-based stage rejects sessions with $< 3$ EDA tool invocations or $> 50\%$ permission-denied lines. A 0.8B LLM classifier (Qwen3.5-0.8B-Q8\_0, CPU) performs a binary is\_design\_flow classification. End-to-end cost is $\sim$1.1s/call at 73 tok/s. On the full candidate pool, the prefilter rejects 0\%---confirming that mc $\geq 30$ sessions are almost universally genuine EDA work.
% \textbf{两级预筛器。} 规则级拒绝 EDA 工具调用 $< 3$ 次或权限被拒行 $> 50\%$ 的会话。一个 0.8B LLM 分类器（Qwen3.5-0.8B-Q8\_0，CPU）执行二分类 is\_design\_flow 判定。端到端开销约 1.1 秒/次，73 tok/s。在完整的候选池上，预筛器拒绝率为 0\%——印证了 mc $\geq 30$ 的会话几乎全为真实 EDA 工作。

\subsection{Domain Keyword Filters}
\label{sec:keyword-filters}
% \subsection{领域关键词过滤器}

Table~\ref{tab:keywords} lists the EDA keywords used for session scanning and area-group assignment. The digital-front group uses an expanded keyword set (synthesis, synopsys, tcl, verilog, systemverilog, rtl, sdc, vivado, yosys, openroad) and a relaxed mc threshold ($\geq 15$) to capture digital-front tasks that use cline less intensively.
% 表~\ref{tab:keywords} 列出用于会话扫描和领域组分配的 EDA 关键词。数字前端组使用扩展关键词集（synthesis、synopsys、tcl、verilog、systemverilog、rtl、sdc、vivado、yosys、openroad）并放宽 mc 阈值（$\geq 15$），以捕获使用 cline 不那么密集的数字前端任务。

\begin{table}[ht]
\centering
\small
\caption{EDA keyword filters by area group.}
% \caption{各领域组的 EDA 关键词过滤器。}
\label{tab:keywords}
\begin{tabular}{@{}llp{0.38\linewidth}@{}}
\toprule
\textbf{Area group} & \textbf{mc threshold} & \textbf{Keywords} \\
\midrule
Digital front & $\geq 15$ & verdi, vcs, dc\_shell, innovus, primetime, synthesis, tcl, verilog, rtl \\
SPICE / char.\ & $\geq 30$ & hspice, finesim, siliconsmart, spectre, ngspice, simulator, spice \\
Physical / layout & $\geq 30$ & calibre, virtuoso, aether, ipcustom, layout, lvs, drc \\
Schematic / analog & $\geq 30$ & waveview, cosmos, skill, netlist, spiros, spectrum, analog, schematic \\
\bottomrule
\end{tabular}
\end{table}

\subsection{Full Domain Distribution}
\label{sec:full-distribution}
% \subsection{完整领域分布}

The benchmark is built in four construction shards. As of 2026-11-21 the living
build holds 2043 cases (A: 443, B: 483, C: 564, D: 554; 36 retired, all in shard A, 2007 active) drawn from the
trace pool described in Section~\ref{sec:scan-pipeline}; these per-shard counts
grow as new traces are curated and reflect the snapshot at submission time (the
main paper reports only the qualitative Full-set scale, not this count). The
Lite subset of \LiteSet{} tasks (used for arm-level evaluation,
Section~\ref{sec:experiments}) is filtered from this living build by
domain-coverage balance and chain-length distribution. The evaluation in
Section~\ref{sec:experiments} uses the frozen \Neval{}-run snapshot of
2026-11-09 (Table~\ref{tab:full-domain-results}). Per-shard counts are maintained
in the benchmark repository; a full user/host distribution table is deferred to
the camera-ready version because the host inventory is company-internal.
% 基准由四个构建分片组成。截至 2026-11-21，活集含 2043 例（A: 443、B: 483、C: 564、
% D: 554；其中 36 例 retired 全属 shard A，active 2007），全部取自第~\ref{sec:scan-pipeline}~节描述的轨迹池；各分片计数随新 trace 整理而增长，
% 此处反映投稿时快照（正文仅报告 Full 集定性规模，不报告此计数）。用于臂级评测的 Lite 子集
%（\LiteSet{} 任务，第~\ref{sec:experiments}~节）即从该活集按领域覆盖平衡与链长分布筛选。
% 第~\ref{sec:experiments}~节的评测使用 2026-11-09 冻结的 \Neval{}-run 快照
%（表~\ref{tab:full-domain-results}）。各分片计数在基准仓库中维护；完整的用户/主机分布表推迟到
% camera-ready，因为主机清单属公司内部信息。
% T12 resolved 2026-11-13：正文中所有工作站/用户/会话级 infra 数字已模糊处理，
% 仅保留方法论描述；Full/Lite 集叙事替代具体扫描统计。

\subsection{Task Schema Specification}
\label{sec:task-schema-spec}
% \subsection{任务模式规范}

Each task is a JSONL record with the following fields:
% 每个任务是一条 JSONL 记录，包含以下字段：
\begin{itemize}
    \item \texttt{id}: unique task identifier (area-group-prefixed, e.g., \texttt{A-001}).
    % \texttt{id}：唯一任务标识（领域组前缀，如 \texttt{A-001}）。
    \item \texttt{mode}: always \texttt{design\_flow}.
    % \texttt{mode}：恒为 \texttt{design\_flow}。
    \item \texttt{prompt}: natural-language task description.
    % \texttt{prompt}：自然语言任务描述。
    \item \texttt{assets}: list of input file paths (copied into the workdir).
    % \texttt{assets}：输入文件路径列表（拷入工作目录）。
    \item \texttt{entry\_point}: \texttt{outputs/submit.zip}.
    % \texttt{entry\_point}：\texttt{outputs/submit.zip}。
    \item \texttt{test}: \texttt{check(deliverable, workdir)} function specification.
    % \texttt{test}：\texttt{check(deliverable, workdir)} 函数规格。
    \item \texttt{test\_rationale}: human-readable assertion list in ``k assertions: (1), (2), ...'' format ($k \geq 8$).
    % \texttt{test\_rationale}：可读断言列表，格式为"k 条断言：(1)、(2)、……"（$k \geq 8$）。
    \item \texttt{trace\_source}: \{host, user, session\_file, session\_msgs, real\_task\_summary\}.
    % \texttt{trace\_source}：\{host, user, session\_file, session\_msgs, real\_task\_summary\}。
    \item \texttt{domain}: one of 13 EDA domains (Section~\ref{sec:full-distribution}).
    % \texttt{domain}：13 个 EDA 领域之一（第~\ref{sec:full-distribution}~节）。
\end{itemize}
